\section{Discussion}\label{sec:discussion}

\textbf{What the information bound changes.} The formulation in terms of $P(N)$ and $N_{0.9}$ makes the statement ``the method finds a hidden attractor'' measurable, and \cref{thm:info} turns it into a question about information: a method can beat uniform sampling only to the extent that its observations distinguish points around the target. Three consequences follow. Adaptivity without auxiliary information is worth at most a factor of 2.6, so a search heuristic can gain only to the extent that what it observes varies near the target; evolutionary search with Zelinka's cost function~\cite{Zelinka2016}, whose cost does not point toward the rare targets of the benchmark, is no better than uniform sampling (\cref{tab:gains}). Uniform sampling is a strong reference: its cost $2.3/p$ is known in advance and does not depend on the geometry of the basin. Finally, the value of a channel can be bounded before any search from a uniform census, as for escape times through $\mu(Q_\infty)$. Rarity, moreover, is a property of the chosen region rather than of the attractor: enlarging the region or moving the parameter toward the crisis of the attractor lowers the basin measure by at most about two orders of magnitude (Section~\ref{app:enlarged}), while a truly rare hidden target arises from a combination of hidden regimes, as in coupled units.

\textbf{Where adaptive search gains by orders of magnitude.} \Cref{prop:product} identifies this class: for $k$ units, the information in escape times is worth $\Theta(q^{-k})$, and ideal splitting attains it, so the bound of \cref{thm:info} is sharp up to a constant; the measured gains grow with base 14 against the ideal 21. The base $1/q$ is a property of a single unit that can be measured before any search, so the proposition is predictive in a way that the diagnostics of \cref{sec:diagnostics}, which need a known target, are not. For single units, the bound of \cref{cor:channels} is moderate (44--230 on the escaping systems of the benchmark), and the measured gains of 3--7 reflect this. Channel B gives an order-of-growth gain where the transient duration grows on all scales as the basin is approached (D2), but the gain is limited in $p$ by the time resolution of the coordinate. Where trajectories are bounded, escape times are worthless, and the settling part of the coordinate is what remains: it gives 7.4 on Chua~II and 55--360 on systems of two and three Chua circuits, where the escape step is the same for every trajectory. Channel A works if the boundary of the known basins is short and the target forms a noticeable part of it, and channel C helps when the heuristic is fortunate (nullclines for Chua~I) and costs nothing when it is not.

\textbf{Why a portfolio.} Single adaptive strategies lose to uniform sampling by a factor of several in the worst cases and produce heavy tails (\cref{sec:failures}). By \cref{prop:portfolio}, the portfolio is guaranteed to cost at most $1/\theta$ times as much as uniform sampling in integrations and inherits the order of every autonomous arm, and \cref{cor:sas} extends this to stagger-and-step chains, which are as strong as splitting or stronger where the signal is present. What matters for the gain is the coordinate; the portfolio removes the risk associated with any particular way of climbing it. Bisection toward basin boundaries~\cite{NusseYorke1989,Skufca2006Edge} is a strong baseline on systems with smooth boundaries but does not scale to coupled units (Section~\ref{app:ablation}).

\textbf{When the gain is a gain in time.} Where trajectories are bounded, as in dissipative physical systems, a nontarget integration is as expensive as a target one, and the gain carries over to time: for the rare regime of two coupled Chua circuits, it is 10 to 68, depending on the evaluation scheme of the reference and on the number of threads. Counting integrations overstates the advantage where most trajectories escape quickly, because the method selects the long ones by construction; there, the gain in integrations becomes a draw or a loss in time unless it is very large. The search strategy and the evaluation scheme are independent choices: the portfolio needs only a regime label and a transient duration for each point, which a recurrence-based mapper also provides, so the two can be combined.

\textbf{Relation to continuation.} The established way to find a rare attractor is to continue it from parameter values where it is easy to find~\cite{Brzeski2017,Datseris2023Framework,Datseris2026Global}. Our setting is that of a fixed parameter value, as when a new system is explored or when the attractor exists only in a narrow parameter window. The approaches are complementary: attractors found at neighboring parameter values are natural prior sets, and the portfolio can replace the uniform sample within each continuation step. For coupled units whose attractors are known, product initial conditions are another prior set; in the coupled Chua circuits, they reach both the rare family and the coupling-born attractor directly. The portfolio is intended for the case in which the units' attractors are not known or the regime is not a product.

\textbf{Verification.} Without a verification protocol, a comparison of search methods loses its meaning: false hidden attractors due to transient chaos are known~\cite{SprottMunmuangsaen2018Comment,Danca2016Transient}, and false absences due to the integration step are, according to our data, just as real (Section~\ref{app:verification}). The same protocol, applied to enumerate the regimes of a new system, found the coupling-born attractor of \cref{sec:newattr}.

\textbf{Limitations.} The conclusions about gains refer to the benchmark, the synthetic families, and the coupled copies; on other systems, only the worst-case bound of \cref{prop:portfolio} is guaranteed, and only in integrations. \Cref{prop:product} assumes exact conditional sampling and uncoupled units; the measured gains follow its geometric law with a smaller base, and in our examples, weak coupling changes the basin measure by less than 30\%. For maps and for systems with escape by slow drift, the reaction coordinate may be uninformative or deceptive. The largest gains are obtained on weakly coupled or uncoupled copies, whose rare target is close to a product of the units' attractors; declaring product states known improves the gains further, by factors of 1.2--1.7 where both variants were measured, and the numbers obtained without this knowledge are reported alongside wherever they were measured. The same formulation applies to trained models of dynamics, such as reservoir computers in closed loop~\cite{Pathak2017Chaos,LuHuntOtt2018}, which may develop attractors absent from the training data~\cite{HaluszczynskiRath2019,HerteuxRath2020,Flynn2021,Kabayama2025PseudoLorenz}; that application will be reported separately.
